%------------------------------------------------------------------------------%
\begin{question}
  Determine a posição relativa e o ponto de interseção, caso exista,
  \[
    r\colon
    \begin{cases}
      x = 1 + t  \\
      y = 2 +2t
    \end{cases}
  \]
  \[
    s\colon
    \begin{cases}
      x = 4 - \tau  \\
      y = 8 -3\tau
    \end{cases}
  \]
\end{question}

%------------------------------------------------------------------------------%
\begin{resolution}{Solução}
  Os vetores diretores são
  \[
    \vec v=(1,2)
  \]
  \[
    \vec w=(-1,-3)
  \]

  \pause
  Eles não são paralelos

  \pause
  Portanto, no plano, as retas são concorrentes
\end{resolution}

%------------------------------------------------------------------------------%
\begin{resolution}{Ponto de Interseção}
\begin{columns}
\column{0.5\textwidth}
  \[\onslide<+->{
    \begin{cases}
      1 +  t = 4 -  \tau \\
      2 + 2t = 8 - 3\tau
    \end{cases}}
  \]
  \[\onslide<+->{
    \begin{cases}
       t = 3 - \tau \\
      2t + 3\tau = 6
    \end{cases}}
  \]
  \begin{align*}
    \onslide<+->{2(3 - \tau) + 3\tau & = 6 \\}
    \onslide<+->{6 - 2\tau + 3\tau   & = 6 \\}
    \onslide<+->{\tau                & = 0}
  \end{align*}
\column{0.5\textwidth}
  \[
   \onslide<+->{t = 3 - \tau}
   \onslide<+->{= 3}
  \]
  \onslide<+->{Substituindo na reta $s$}
  \begin{align*}
    \onslide<+->{x & = 4 - \tau}\onslide<+->{ = 4} \\
    \onslide<+->{y & = 8 -3\tau}\onslide<+->{ = 8}
  \end{align*}
  \onslide<+->{Assim
  \(
    r\cap s = \{(4, 8)\}
  \)}
\end{columns}
\end{resolution}

%------------------------------------------------------------------------------%
\endinput
